\begin{abstract}
As Taiwan enters a super-aged society, the real burden of curbside waste collection is not locating the garbage truck but deciding whether a resident, given their location and walking pace, can still leave home and arrive comfortably in time. This departure decision faces three compounding uncertainties: arrival-time prediction, walking time, and data freshness. Rather than refining arrival-time accuracy, this study defines when a system is entitled to recommend departure and when it must degrade. Using Taichung City as the study area, we build a domain-specific artificial intelligence agent that encapsulates waste-collection knowledge, including collection schedules, pedestrian network accessibility, and operational characteristics, as whitelisted tools. A large language model handles only intent recognition and explanation, whereas numerical computation and decisions are executed by controlled tools with schema-validated inputs and outputs, confining tool hallucination within verifiable boundaries. Feasibility is judged by a time margin, namely the predicted arrival time minus the walking time and a safety buffer, with the buffer adapted to data latency and prediction error; walking time is estimated with a spatial database. Rolling-origin backtesting compares historical-average, distance--speed, and regression predictors by mean absolute error and root-mean-square error, treats the false-feasibility rate as the primary safety metric, reports recommendation coverage and fallback success rate, and includes at least thirty anomaly scenarios. Expected outcomes include statistically significant improvement over the historical-average baseline, a bounded false-feasibility rate without collapse in coverage, full disclosure of data sources and fallback reasons, and an explicit applicability boundary for personalized departure advice under low-frequency data.
\end{abstract}
\begin{IEEEkeywords}
Arrival time prediction, decision support systems, domain-specific agents, graceful degradation, large language models, tool hallucination, waste collection.
\end{IEEEkeywords}
